\section{Introducción}\label{sec:introduccion}

El diseño de algoritmos tiene como objetivo encontrar soluciones a problemas de manera eficiente teniendo en cuenta las limitaciones de tiempo y memoria que existen dentro de los sistemas informáticos.

Tres de los problemas más comunes en la computación son el \textbf{ordenamiento} de datos, la \textbf{búsqueda} de información y la \textbf{selección} de elementos dentro de estructuras de datos; estos problemas se pueden afrontar desde distintas perspectivas, y en este informe los abordamos desde el paradigma de \textbf{Divide y Vencerás}.

El presente informe aborda estos problemas sobre una base de datos ficticia de deportistas (generada de manera procedimental) con el fin de exponer, analizar y comparar algoritmos de ordenamiento (\textit{Merge sort} clásico y optimizado, \textit{Quick sort} con múltiples estrategias de pivote), algoritmos de búsqueda (\textit{Búsqueda binaria recursiva}, \textit{exponencial} e \textit{interpolación}) y selección (\textit{Quick select}). Todo ello analizando su comportamiento a nivel teórico y contrastándolo mediante experimentación empírica.
